% Provenance (never printed): see sources/outline_sources.md (Year 2: radiative forcing
% quantified; feedbacks and climate sensitivity; detection and attribution; scenarios; the
% remaining carbon budget; components of sea-level rise; physical science only). Every statement
% from IPCC AR6 WG1 (SPM; chapters 3, 5, 7 and 9; Annex III data), ledger rows in
% sources/ledger/book2.md (ipcc:*) and sources/ledger/book4.md (ar6*:*). Curves and numbers from
% figdata/bachelor-2/anthropogenic.py and the Book 4 second-writer scratch calculator calc/ch27.py.
\chapter{Antropogene klimaatverandering}\label{ch:b2:anthropogenic-climate-change}

Sinds 1970 is het aardoppervlak sneller opgewarmd dan in enige andere periode van vijftig jaar in minstens
de laatste tweeduizend jaar (\emph{hoog vertrouwen}). De schoolboeken vertelden wat er is veranderd en
waarom. Dit hoofdstuk maakt van dat verhaal een balans: hoeveel watt per vierkante meter de menselijke
drijvende factoren toevoegen, hoeveel de warmere planeet terugstuurt naar de ruimte, hoeveel opwarming elke ton
koolstofdioxide oplevert, en hoe de stijgende zee wordt verdeeld tussen opwarmend water en smeltend ijs. Elk getal
komt met zijn onzekerheid, in de gekalibreerde taal van de beoordelingen.
% ledger: ipcc:fastest-2000yr, ipcc:likelihood

\begin{recall}
\textbf{Uit de schoolboeken}: stralingsforcering, $\qty{3.93}{W.m^{-2}}$ voor een verdubbeling van
koolstofdioxide, en haar logaritmische wet; klimaatterugkoppeling; \omterm{def:b2:anthropogenic-climate-change:ecs}{evenwichtsklimaatgevoeligheid}, beste
schatting \qty{3}{\celsius}, \emph{waarschijnlijk} 2.5 tot \qty{4}{\celsius}; \omterm{def:b2:anthropogenic-climate-change:da}{attributie}, emissiescenario
en projectie; de gekalibreerde taal: \emph{waarschijnlijk} 66 tot \qty{100}{\percent}, \emph{zeer
waarschijnlijk} 90 tot \qty{100}{\percent}, waarbij een bereik $[x \text{ tot } y]$ een zeer waarschijnlijk bereik is. De menselijke
invloed heeft het klimaat tussen 1850--1900 en 2010--2019 met \qty{1.07}{\celsius} opgewarmd.
\textbf{Uit de natuurkundeboeken}: het zwarte lichaam en de wet van Stefan--Boltzmann; thermische uitzetting.
\textbf{Uit de biologieboeken}: de koolstofcyclus, de luchtfractie, koolstofputten en het
koolstofbudget van de biosfeer. \textbf{Uit \cref{ch:b2:deep-time-climate}}: de trage koolstofcyclus,
die een overmaat koolstofdioxide pas over honderdduizenden jaren terugneemt.
\end{recall}
% ledger: ipcc:erf-2xco2, ipcc:log-co2, ipcc:ecs, ipcc:likelihood, ipcc:attribution

\section{De forcering gekwantificeerd}\label{sec:b2:anthropogenic-climate-change:forcing}

\begin{definition}[Effectieve stralingsforcering]\label{def:b2:anthropogenic-climate-change:erf}
De \emph{effectieve stralingsforcering}\index{effectieve stralingsforcering} (ERF) van een verstoring
--- een verandering van een broeikasgas, van aerosolen, van de zon --- is de verandering die ze veroorzaakt in de netto
neerwaartse stralingsflux aan de top van de atmosfeer, in \unit{W.m^{-2}}, zodra de atmosfeer,
de wolken en het landoppervlak zich eraan hebben aangepast, maar voordat de globale oppervlaktetemperatuur
heeft gereageerd. De aanpassingen die worden meegeteld, zijn die welke door de factor zelf worden veroorzaakt, onafhankelijk van de
oppervlaktetemperatuur.
\end{definition}
% ledger: ar6ch7:erf-def

De definitie trekt de grens tussen forcering en respons volgens wat van de
oppervlaktetemperatuur afhangt, niet volgens de snelheid. De ogenblikkelijke forcering van koolstofdioxide is de verandering van de flux net
nadat het gas is toegevoegd; de stratosfeer, die dan binnen enkele maanden afkoelt, verandert ze een eerste keer;
de troposfeer, haar waterdamp en haar wolken veranderen ze opnieuw. Voor een verdubbeling van
koolstofdioxide vindt de beoordeling \qty{3.75}{W.m^{-2}} zodra de stratosfeer zich heeft aangepast, en
aanpassingen in de troposfeer die elkaar deels opheffen: \qty{-0.60}{W.m^{-2}} door haar opwarming, +0.22
door waterdamp, +0.45 door wolken en +0.11 door het landoppervlak, samen \qty{+0.18}{W.m^{-2}}.
De ERF is $3.75 + 0.18 = \qty{3.93}{W.m^{-2}}$, \qty{6}{\percent} meer dan de
\qty{3.71}{W.m^{-2}} van de vorige beoordeling (\cref{exo:b2:anthropogenic-climate-change:1}).
% ledger: ar6ch7:2xco2, ar6ch7:erf-def

De forcering van elke factor sinds 1750 is jaar voor jaar gereconstrueerd. De figuur hieronder toont
haar: koolstofdioxide op \qty{2.16}{W.m^{-2}} in 2019, de andere broeikasgassen --- methaan,
distikstofoxide, de gehalogeneerde gassen en ozon --- samen op \qty{1.69}{W.m^{-2}}, en de
aerosolen op \qty{-1.06}{W.m^{-2}}. De totale menselijke forcering is \qty{2.72}{W.m^{-2}}, en haar zeer
waarschijnlijke bereik, 1.96 tot \qty{3.48}{W.m^{-2}}, is breed vooral door de aerosolen. Hun ERF over
1750--2014 wordt beoordeeld op $-1.3$ [$-2.0$ tot $-0.6$] \unit{W.m^{-2}} (\emph{middelgroot vertrouwen}),
grotendeels via hun effect op wolken, die helderder worden wanneer aerosolen meer druppeltjes leveren.
Vulkaanuitbarstingen voegen korte negatieve pieken toe (\cref{ch:b2:eruptive-styles}): de diepste in de
reeks, \qty{-4.8}{W.m^{-2}} in 1816, volgt op de uitbarsting van de Tambora in 1815.
% ledger: data:b4-ar6-erf, ipcc:erf-total, ipcc:erf-agents, ar6ch7:aerosol, usgs:msh-comparisons

\begin{omfigure}
\centering
\begin{tikzpicture}
\begin{axis}[omrecord, width=0.92\linewidth, height=6.2cm, xmin=1750, xmax=2020, ymin=-5.2,
  ymax=3.8, xlabel={jaar}, ylabel={ERF sinds 1750 (\unit{W.m^{-2}})},
  xtick={1750,1800,...,2000}, /pgf/number format/1000 sep={}, xtick pos=bottom, ytick pos=left,
  tick label style={font=\scriptsize}, label style={font=\small},
  legend style={at={(0.5,-0.24)}, anchor=north, font=\scriptsize, draw=black!30,
  legend columns=3, /tikz/every even column/.append style={column sep=6pt}},
  legend cell align=left]
% computed in figdata/bachelor-2/anthropogenic.py part erf; ledger: data:b4-ar6-erf
\addplot[name path=hi, draw=none, forget plot] table[x=year, y=hi] {figdata/out/bachelor-2-anthropogenic-erf.dat};
\addplot[name path=lo, draw=none, forget plot] table[x=year, y=lo] {figdata/out/bachelor-2-anthropogenic-erf.dat};
\addplot[black!15, forget plot] fill between[of=hi and lo];
\addplot[black!45, line width=0.4pt] table[x=year, y=natural] {figdata/out/bachelor-2-anthropogenic-erf.dat};
\addlegendentry{natuurlijk (vulkanen, zon)}
\addplot[omThm, line width=1pt] table[x=year, y=co2] {figdata/out/bachelor-2-anthropogenic-erf.dat};
\addlegendentry{koolstofdioxide}
\addplot[omThm!50!orange, line width=1pt, densely dashed] table[x=year, y=other_ghg] {figdata/out/bachelor-2-anthropogenic-erf.dat};
\addlegendentry{andere broeikasgassen en ozon}
\addplot[omDef, line width=1pt] table[x=year, y=aerosol] {figdata/out/bachelor-2-anthropogenic-erf.dat};
\addlegendentry{aerosolen}
\addplot[black, line width=1.4pt] table[x=year, y=anthro] {figdata/out/bachelor-2-anthropogenic-erf.dat};
\addlegendentry{totaal menselijk (grijs: 5--95\,\%)}
\draw[black!40] (axis cs:1750,0) -- (axis cs:2020,0);
\node[anchor=west, font=\scriptsize, text=black!60, inner sep=1pt] at (axis cs:1817.5,-4.75) {$\leftarrow$ de Tambora};
\end{axis}
\end{tikzpicture}
\par\smallskip{\small \omterm{def:b2:anthropogenic-climate-change:erf}{Effectieve stralingsforcering} van 1750 tot 2019, ten opzichte van 1750, per factor.
De totale menselijke forcering bereikt \qty{2.72}{W.m^{-2}} in 2019, met het zeer waarschijnlijke bereik gearceerd.
Gegevens: IPCC Sixth Assessment Report, Working Group~I, Annex~III extended data (CC BY 4.0).}
\end{omfigure}
% values: anthro 2.7204 (1.9569-3.4819), co2 2.1563, other_ghg 1.6850, aerosol -1.0583, volcanic -4.8162 in 1816 (anthropogenic.py numbers)

\section{De energiebalans}\label{sec:b2:anthropogenic-climate-change:energy}

Een positieve forcering laat de planeet energie winnen. Terwijl ze opwarmt, straalt ze meer uit naar de ruimte, en de
winst krimpt; het klimaat komt pas tot rust wanneer de twee in evenwicht zijn.

\begin{definition}[Energieonbalans van de aarde; klimaatterugkoppelingsparameter]\label{def:b2:anthropogenic-climate-change:eei}
De \emph{energieonbalans van de aarde}\index{energieonbalans van de aarde} $N$ is de netto energieflux het
klimaatsysteem in aan de top van de atmosfeer, gemiddeld over de aarde en over een jaar of langer,
in \unit{W.m^{-2}}: geabsorbeerd zonlicht min uitgaande infrarode straling. Ze is de snelheid waarmee de oceaan,
het land, het ijs en de lucht warmte winnen, per vierkante meter aardoppervlak.
\end{definition}

\begin{definition}[Klimaatterugkoppelingsparameter]\label{def:b2:anthropogenic-climate-change:lambda}
De \emph{klimaatterugkoppelingsparameter}\index{klimaatterugkoppelingsparameter} $\lambda$ is de verandering van
de netto neerwaartse flux aan de top van de atmosfeer per graad globale oppervlakteopwarming, in
\unit{W.m^{-2}.\celsius^{-1}}, veroorzaakt door de respons van het klimaatsysteem op die opwarming. Een
stabiel klimaat heeft $\lambda < 0$: opwarming vergroot het energieverlies naar de ruimte.
\end{definition}

Het schoolboek schreef de balans met de grootte van $\lambda$, een positief getal; vanaf hier wordt het
teken behouden, zoals in de beoordelingen, die de parameter $\alpha$ noemen.
% ledger: ar6ch7:energy-budget

\begin{proposition}[Globale energiebalans]\label{prop:b2:anthropogenic-climate-change:balance}
Laat $\Delta F$ de \omterm{def:b2:anthropogenic-climate-change:erf}{effectieve stralingsforcering} zijn en $\Delta T$ de globale oppervlakteopwarming
ten opzichte van een evenwichtstoestand. Als de respons van de uitgaande flux lineair is in $\Delta T$, dan geldt
\[
  \Delta N = \Delta F + \lambda\,\Delta T .
\]
In een nieuw evenwicht is $\Delta N = 0$, dus $\Delta T_{\text{eq}} = -\Delta F/\lambda$.
\end{proposition}
\begin{proof}
Per definitie verandert de forcering de netto neerwaartse flux met $\Delta F$ bij vaste
oppervlaktetemperatuur, en de respons op de opwarming verandert hem met $\lambda\,\Delta T$ in eerste orde;
de onbalans is hun som. Ze gelijkstellen aan nul geeft het evenwicht.
\end{proof}

$N$ meten betekent de warmte meten die het systeem wint, en de oceaan bevat er bijna alles van.
Van 1971 tot 2018 won het klimaatsysteem \qty{435}{ZJ} (zettajoule, \qty{e21}{J}): de oceaan
\qty{91}{\percent}, \qty{242}{ZJ} in zijn bovenste \qty{700}{m} en \qty{154}{ZJ} daaronder; het land
\qty{5}{\percent}, het smeltende ijs \qty{3}{\percent}, de atmosfeer \qty{1}{\percent}. De gemiddelde
onbalans bedroeg \qty{0.57}{W.m^{-2}} over 1971--2018, en ze is gegroeid: \qty{0.79}{W.m^{-2}}
[0.52 tot 1.06] over 2006--2018 (\cref{exo:b2:anthropogenic-climate-change:2}).
% ledger: ar6ch7:inventory, ipcc:heating

\begin{omfigure}
\centering
\begin{tikzpicture}[font=\small, >=Stealth]
\draw[fill=omDef!8, draw=omDef!60] (-3.2,0) rectangle (3.2,1.6);
\node[align=center] at (0,1.05) {klimaatsysteem};
\node[align=center, font=\scriptsize] at (0,0.5) {oceaan \qty{91}{\percent}, land \qty{5}{\percent}, ijs \qty{3}{\percent}, lucht \qty{1}{\percent}};
\draw[->, line width=2.2pt, omThm] (-2.0,3.3) -- node[pos=0.12, left=3pt, align=right, font=\scriptsize] {forcering $\Delta F$} (-2.0,1.7);
\draw[->, line width=2.2pt, omDef] (2.0,1.7) -- node[pos=0.88, right=3pt, align=left, font=\scriptsize] {respons $\lambda\,\Delta T$} (2.0,3.3);
\draw[->, line width=1.4pt, black!70] (0,-0.1) -- node[right=3pt, align=left, font=\scriptsize] {opgeslagen: $\Delta N = \Delta F + \lambda\,\Delta T$} (0,-1.0);
\draw[black!40, dashed] (-3.4,2.4) -- (3.4,2.4) node[right, font=\scriptsize, text=black!60] {top van de atmosfeer};
\end{tikzpicture}
\par\smallskip{\small De energiebalans van een geforceerd klimaat (schematisch). De forcering voegt energie toe aan
de top van de atmosfeer, de respons van de warmere planeet haalt er een deel weg ($\lambda\,\Delta T$ is
negatief), en het verschil wordt opgeslagen, grotendeels in de oceaan. Aandelen van de gewonnen warmte van 1971
tot 2018.}
\end{omfigure}
% ledger: ar6ch7:inventory

\begin{method}[Een klimaatgevoeligheid uit de historische energiebalans]\label{met:b2:anthropogenic-climate-change:energy-budget}
\begin{enumerate}
  \item Kies twee perioden, een referentie dicht bij evenwicht en een recente; neem de verandering van de
        forcering $\Delta F$, van de onbalans $\Delta N$ en van de globale temperatuur $\Delta T$.
  \item Los de balans op naar de terugkoppelingsparameter: $\lambda = (\Delta N - \Delta F)/\Delta T$.
  \item De gevoeligheid voor een verdubbeling is dan $-F_{2\times}/\lambda$; de opwarming op het moment van de
        forcering, geschaald naar $F_{2\times}$, is $F_{2\times}\Delta T/\Delta F$.
\end{enumerate}
\end{method}

Van 1850--1900 tot 2006--2019 neemt de beoordeling $\Delta N = \qty{0.59}{W.m^{-2}}$, $\Delta T =
\qty{1.03}{\celsius}$ en $\Delta F = \qty{2.20}{W.m^{-2}}$. Dan is $\lambda = (0.59 - 2.20)/1.03 =
\qty{-1.56}{W.m^{-2}.\celsius^{-1}}$, en de gevoeligheid is $3.93/1.56 = \qty{2.5}{\celsius}$,
de waarde die de beoordeling vindt, met een zeer waarschijnlijk bereik van 1.6 tot \qty{4.8}{\celsius}. Maar dit is
een \emph{effectieve} gevoeligheid, voor de terugkoppelingen van de opwarming tot nu toe. De opwarming was
het sterkst waar de terugkoppelingen het meest negatief zijn, in de westelijke tropische Stille Oceaan, en zwakker boven de
oostelijke Stille Oceaan en de Zuidelijke Oceaan; naarmate dat patroon verandert, wordt verwacht dat de terugkoppelingsparameter
minder negatief wordt (\emph{hoog vertrouwen}), dus de gevoeligheid op lange termijn is hoger dan de
historische schatting (\cref{exo:b2:anthropogenic-climate-change:7}).
% ledger: ar6ch7:energy-budget, ar6ch7:pattern, ipcc:erf-2xco2
% values: lambda_energy -1.563, ecs_energy 2.514 (anthropogenic.py numbers)

\section{Terugkoppelingen en gevoeligheid}\label{sec:b2:anthropogenic-climate-change:feedbacks}

\begin{definition}[Planck-respons; terugkoppelingen van waterdamp, temperatuurgradiënt en wolken]\label{def:b2:anthropogenic-climate-change:feedbacks}
De terugkoppelingsparameter wordt opgesplitst in de bijdragen van de grootheden die met de
oppervlaktetemperatuur veranderen. De \emph{Planck-respons}\index{Planck-respons} is de extra infrarode straling die
naar de ruimte wordt uitgezonden wanneer het oppervlak en de hele troposfeer evenveel opwarmen. De
\emph{waterdampterugkoppeling}\index{waterdampterugkoppeling} is de verandering van de flux door de toename
van de atmosferische waterdamp met de temperatuur. De \emph{terugkoppeling van de temperatuurgradiënt}\index{terugkoppeling van de temperatuurgradiënt}
is de verandering van de flux door een opwarming die niet gelijkmatig is met de hoogte. De
\emph{wolkenterugkoppeling}\index{wolkenterugkoppeling} is de verandering van de flux door de verandering van de hoeveelheid, de hoogte en
de helderheid van de wolken.
\end{definition}

\begin{proposition}[De Planck-respons]\label{prop:b2:anthropogenic-climate-change:planck}
Als de aarde een vaste fractie $1 - g$ van de infrarode straling $\sigma T^4$ van haar oppervlak bij
temperatuur $T$ naar de ruimte uitzendt, is haar \omterm{def:b2:anthropogenic-climate-change:feedbacks}{Planck-respons}
\[
  \lambda_{\text P} = -(1 - g)\,4\sigma T^3 .
\]
Met $g = 0.4$ en $T = \qty{288}{K}$ is $\lambda_{\text P} = \qty{-3.25}{W.m^{-2}.K^{-1}}$.
\end{proposition}
\begin{proof}
De uitgaande flux is $(1 - g)\sigma T^4$; de afgeleide ervan naar $T$ is $(1 - g)\,4\sigma
T^3$, en een extra verlies is een negatieve bijdrage aan de netto neerwaartse flux. Numeriek is $0.6
\times 4 \times \num{5.670e-8} \times 288^3 = 3.25$.
\end{proof}
% ledger: ar6ch7:planck, const:sigma, ipcc:budget-surface
% values: planck -3.2509 (anthropogenic.py numbers)

De fractie $g$ is het genormaliseerde broeikaseffect: de atmosfeer houdt \qty{159}{W.m^{-2}} vast van
de \qty{398}{W.m^{-2}} die het oppervlak uitzendt, $g \approx 0.4$. De beoordeling vindt met gedetailleerde
stralingscodes $-3.22$ [$-3.4$ tot $-3.0$] (\emph{hoog vertrouwen}). Waterdamp volgt
op mondiale schaal de wet van Clausius--Clapeyron, en meer damp absorbeert meer infrarood: de
\omterm{def:b2:anthropogenic-climate-change:feedbacks}{waterdampterugkoppeling} is sterk positief, \qty{1.85}{W.m^{-2}.\celsius^{-1}} volgens satellieten.
In de tropen warmt de bovenste troposfeer meer op dan het oppervlak en straalt daardoor meer uit: de
\omterm{def:b2:anthropogenic-climate-change:feedbacks}{terugkoppeling van de temperatuurgradiënt} is daar negatief, positief dicht bij de polen, en globaal ongeveer $-0.50$. De twee
worden samen beoordeeld, \num{1.30} [1.1 tot 1.5], omdat wat een model in de ene stopt, het vaak
uit de andere haalt. Smeltende sneeuw en ijs maken het oppervlak donkerder, +0.35; wolken voegen \num{0.42}
toe [$-0.10$ tot $0.94$], de grootste onzekerheid. Een netto negatieve \omterm{def:b2:anthropogenic-climate-change:feedbacks}{wolkenterugkoppeling} is \emph{zeer onwaarschijnlijk}.
% ledger: ar6ch7:planck, ar6ch7:wv-lr, ar6ch7:table710, ipcc:cloud-negative, ipcc:budget-surface

\begin{omfigure}
\centering
\begin{tikzpicture}
\begin{axis}[xbar, width=0.78\linewidth, height=5.4cm, xmin=-3.7, xmax=1.8, ymin=0.4, ymax=6.6,
  ytick={1,...,6}, yticklabels={netto, biogeofysisch en chemisch, wolken, albedo van het oppervlak,
  waterdamp + temperatuurgradiënt, Planck}, y tick label style={font=\scriptsize},
  tick label style={font=\scriptsize}, xlabel={terugkoppelingsparameter (\unit{W.m^{-2}.\celsius^{-1}})},
  label style={font=\small}, bar width=7pt, axis x line=bottom, axis y line=left,
  axis line style={-}, xmajorgrids, grid style={gray!22, line width=0.3pt}, extra x ticks={0},
  extra x tick style={grid style={black!60, line width=0.5pt}}]
% ledger: ar6ch7:table710 (central values and very likely ranges)
\addplot[fill=omDef!55, draw=omDef!80!black, error bars/.cd, x dir=both, x explicit,
  error bar style={black!70}] coordinates {
  (-3.22,6) += (0.22,0) -= (0.18,0)
  (-0.01,2) += (0.26,0) -= (0.26,0)
  (-1.16,1) += (0.65,0) -= (0.65,0)};
\addplot[fill=omThm!55, draw=omThm!80!black, error bars/.cd, x dir=both, x explicit,
  error bar style={black!70}] coordinates {
  (1.30,5) += (0.20,0) -= (0.20,0)
  (0.35,4) += (0.25,0) -= (0.25,0)
  (0.42,3) += (0.52,0) -= (0.52,0)};
\end{axis}
\end{tikzpicture}
\par\smallskip{\small De klimaatterugkoppelingen volgens de beoordeling van het IPCC (Sixth Assessment Report, Working
Group~I, tabel 7.10), centrale waarden en zeer waarschijnlijke bereiken. Positieve terugkoppelingen versterken een opwarming
(rood), negatieve dempen haar (blauw); de netto waarde is de som.}
\end{omfigure}

\begin{definition}[Evenwichtsklimaatgevoeligheid]\label{def:b2:anthropogenic-climate-change:ecs}
De \emph{evenwichtsklimaatgevoeligheid}\index{evenwichtsklimaatgevoeligheid} (ECS) is de
globale oppervlakteopwarming, in het nieuwe evenwicht, veroorzaakt door een verdubbeling van de koolstofdioxide van de
lucht ten opzichte van het pre-industriële niveau, waarbij de ijskappen vast worden gehouden.
\end{definition}

\begin{proposition}[Gevoeligheid en terugkoppeling]\label{prop:b2:anthropogenic-climate-change:ecs}
Als de terugkoppelingsparameter constant is, geldt $\text{ECS} = -F_{2\times}/\lambda$. Met $F_{2\times} =
\qty{3.93}{W.m^{-2}}$ en de netto $\lambda = \qty{-1.16}{W.m^{-2}.\celsius^{-1}}$ van de terugkoppelingen
hierboven is $\text{ECS} = \qty{3.4}{\celsius}$; de \omterm{def:b2:anthropogenic-climate-change:feedbacks}{Planck-respons} alleen zou $3.93/3.22 =
\qty{1.2}{\celsius}$ geven.
\end{proposition}
\begin{proof}
Pas \cref{prop:b2:anthropogenic-climate-change:balance} toe in evenwicht met $\Delta F =
F_{2\times}$.
\end{proof}
% ledger: ar6ch7:table710, ipcc:erf-2xco2
% values: ecs_process 3.388 (anthropogenic.py numbers)

De terugkoppelingen zijn één bewijslijn; de opwarming sinds 1850, de energiebalans en de klimaten
uit het verleden zijn andere. Samen geven ze de beoordeelde ECS van \qty{3}{\celsius}, \emph{waarschijnlijk}
2.5 tot \qty{4}{\celsius} (\emph{hoog vertrouwen}), \emph{zeer waarschijnlijk} 2 tot \qty{5}{\celsius}.
Omdat de ECS varieert als $1/\lambda$, geeft een symmetrische onzekerheid op $\lambda$ een bereik
van de ECS dat naar hoge waarden is uitgerekt (\cref{exo:b2:anthropogenic-climate-change:12}).
% ledger: ipcc:ecs

\begin{definition}[Transiënte klimaatrespons]\label{def:b2:anthropogenic-climate-change:tcr}
De \emph{transiënte klimaatrespons}\index{transiënte klimaatrespons} (TCR) is de globale
oppervlakteopwarming op het moment dat de koolstofdioxide van de lucht verdubbeld is, wanneer die met \qty{1}{\percent}
per jaar stijgt vanaf het pre-industriële niveau, ongeveer zeventig jaar na het begin.
\end{definition}

De TCR is kleiner dan de ECS omdat de oceaan nog warmte opneemt wanneer de verdubbeling wordt
bereikt. De beste schatting is \qty{1.8}{\celsius}, \emph{waarschijnlijk} 1.4 tot \qty{2.2}{\celsius}
(\emph{hoog vertrouwen}); de historische energiebalans alleen geeft eveneens $3.93 \times 1.03/2.20 =
\qty{1.8}{\celsius}$.
% ledger: ar6ch7:tcr, ar6ch7:energy-budget
% values: tcr_energy 1.840 (anthropogenic.py numbers)

\section{Detectie en attributie}\label{sec:b2:anthropogenic-climate-change:attribution}

\begin{definition}[Detectie; attributie; klimaatvingerafdruk]\label{def:b2:anthropogenic-climate-change:da}
\emph{Detectie}\index{detectie (van klimaatverandering)} van een klimaatverandering is aantonen dat het
klimaat in een welomschreven statistische zin is veranderd, voorbij wat zijn interne variabiliteit zou kunnen
voortbrengen, zonder de oorzaak te geven. \emph{Attributie}\index{attributie} is het beoordelen van de
bijdragen van verscheidene oorzaken aan een gedetecteerde verandering, met een toegekende waarschijnlijkheid. Een
\emph{klimaatvingerafdruk}\index{klimaatvingerafdruk} is het patroon in ruimte en tijd van de respons op één
drijvende factor, zoals gesimuleerd door klimaatmodellen, waaraan die factor in de waarnemingen wordt herkend.
\end{definition}

Verschillende drijvende factoren laten verschillende vingerafdrukken na. Meer broeikasgas warmt de troposfeer op en koelt
de stratosfeer af, omdat een dikkere infrarode deken de bovenste lagen meer naar de ruimte laat uitstralen;
een helderdere zon zou beide opwarmen. De waarnemingen tonen het eerste patroon: de troposfeer is opgewarmd
sinds minstens de jaren 1950, en het is \emph{vrijwel zeker} dat de stratosfeer is afgekoeld. Het is
\emph{zeer waarschijnlijk} dat de menselijke broeikasgassen de belangrijkste oorzaak waren van de opwarming van de troposfeer
sinds 1979, en \emph{extreem waarschijnlijk} dat de afbraak van de ozonlaag de belangrijkste oorzaak was van
de afkoeling van de lagere stratosfeer van 1979 tot het midden van de jaren 1990.
% ledger: ar6ch3:upper-air

\begin{method}[Optimale vingerafdrukmethode]\label{met:b2:anthropogenic-climate-change:fingerprinting}
\begin{enumerate}
  \item Simuleer de respons op elke drijvende factor afzonderlijk: de vingerafdrukken $x_1, \dots, x_k$, als
        vectoren van waarden in ruimte en tijd.
  \item Schrijf de waarnemingen als $y = \sum_i \beta_i x_i + \varepsilon$, met $\varepsilon$ de
        interne variabiliteit, waarvan de statistiek uit lange ongeforceerde simulaties wordt gehaald.
  \item Schat de schaalfactoren $\beta_i$ met de kleinstekwadratenmethode, gewogen met de inverse van de
        variabiliteit zodat de luidruchtigste componenten het minst meetellen, samen met hun betrouwbaarheidsintervallen.
  \item Een $\beta_i$ die significant groter is dan nul, betekent dat de respons op factor $i$ is
        gedetecteerd; een interval dat 1 bevat en 0 uitsluit, betekent dat de gesimuleerde respons
        in overeenstemming is met de waarnemingen, en dat de verandering deels aan die factor kan worden toegeschreven.
\end{enumerate}
\end{method}
% ledger: ar6ch3:fingerprint

Toegepast op de oppervlaktetemperatuur geeft de methode de \omterm{def:b2:anthropogenic-climate-change:da}{attributie} die aan het begin werd herhaald: van de
\qty{1.06}{\celsius} die tussen 1850--1900 en 2010--2019 werd waargenomen, veroorzaakte de menselijke invloed \emph{waarschijnlijk}
0.8 tot \qty{1.3}{\celsius}, beste schatting \qty{1.07}{\celsius}; de natuurlijke drijvende factoren en de
interne variabiliteit dragen niet meer dan enkele tienden van een graad bij, in de ene of de andere richting
(\cref{exo:b2:anthropogenic-climate-change:9,exo:b2:anthropogenic-climate-change:10}).
% ledger: ipcc:attribution, ipcc:observed-2010-2019

\section{Scenario's en cumulatieve uitstoot}\label{sec:b2:anthropogenic-climate-change:scenarios}

\begin{definition}[Gedeeld sociaaleconomisch pad]\label{def:b2:anthropogenic-climate-change:ssp}
Een \emph{gedeeld sociaaleconomisch pad}\index{gedeeld sociaaleconomisch pad} (SSP) is een van vijf
verhalen over hoe bevolking, economieën en technologie zich in de loop van de eeuw zouden kunnen ontwikkelen. De
emissiescenario's van de beoordeling worden SSP$x$-$y$ genoemd: $x$ geeft het pad aan, en $y$ de
benaderde stralingsforcering, in \unit{W.m^{-2}}, die het scenario in 2100 bereikt.
\end{definition}
% ledger: ar6spm:ssp

De vijf belangrijkste scenario's lopen van SSP1-1.9 tot SSP5-8.5. De beoordeling van de fysische wetenschap neemt ze als
gegeven: ze weegt hun waarschijnlijkheid niet en beoordeelt de verhalen niet. Voor elk geven de klimaatmodellen
een geprojecteerde opwarming voor 2081--2100 boven 1850--1900, van \qty{1.4}{\celsius} (zeer waarschijnlijk 1.0
tot 1.8) onder SSP1-1.9 tot \qty{4.4}{\celsius} (3.3 tot 5.7) onder SSP5-8.5. De projecties delen een
eenvoudige wet.
% ledger: ar6spm:ssp, ipcc:proj-temp

\begin{definition}[Transiënte klimaatrespons op cumulatieve CO$_2$-uitstoot]\label{def:b2:anthropogenic-climate-change:tcre}
De \emph{transiënte klimaatrespons op cumulatieve \ce{CO2}-uitstoot}\index{transiënte klimaatrespons op cumulatieve CO$_2$-uitstoot}
(TCRE) is de globale oppervlakteopwarming per eenheid
koolstofdioxide die in totaal sinds de pre-industriële periode is uitgestoten, in \unit{\celsius} per
\qty{1000}{Gt} \ce{CO2}.
\end{definition}

\begin{proposition}[De opwarming volgt de cumulatieve uitstoot]\label{prop:b2:anthropogenic-climate-change:tcre}
In goede benadering is de opwarming door koolstofdioxide evenredig met het totaal dat is uitgestoten,
wat het verloop van de uitstoot ook is: $\Delta T \approx \text{TCRE} \times E$, met $E$ de cumulatieve
uitstoot.
\end{proposition}
\admitted
% ledger: ar6spm:tcre, ar6ch5:tcre-linear

De wet is niet vanzelfsprekend. De forcering groeit slechts met de logaritme van de concentratie, dus elke
toegevoegde ton warmt minder op; maar de oceaan neemt warmte en koolstof steeds minder efficiënt op naarmate de
uitstoot zich ophoopt, zodat meer van elke ton in de lucht blijft en de opwarming ervan minder wordt vertraagd. De
twee effecten heffen elkaar bijna op, en alle studies zijn het erover eens dat dat de oorsprong van de wet is. De betrekking is
bijna lineair (\emph{hoog vertrouwen}), en elke \qty{1000}{Gt} koolstofdioxide veroorzaakt \emph{waarschijnlijk}
0.27 tot \qty{0.63}{\celsius} opwarming, beste schatting \qty{0.45}{\celsius}, of \qty{1.65}{\celsius}
per \qty{1000}{PgC} koolstof (\cref{exo:b2:anthropogenic-climate-change:4}). Van 1850 tot 2019
werd $\num{2390} \pm \qty{240}{Gt}$ koolstofdioxide uitgestoten.
% ledger: ar6ch5:tcre-linear, ar6spm:tcre, ar6spm:budget

\begin{definition}[Resterend koolstofbudget]\label{def:b2:anthropogenic-climate-change:budget}
Het \emph{resterende koolstofbudget}\index{resterend koolstofbudget} voor een temperatuurgrens is de
cumulatieve koolstofdioxide die vanaf een opgegeven datum nog kan worden uitgestoten terwijl de globale opwarming
met een opgegeven waarschijnlijkheid onder die grens blijft, waarbij de opwarming door de andere menselijke drijvende factoren
in rekening wordt gebracht.
\end{definition}

\begin{method}[Het resterende koolstofbudget uit de TCRE]\label{met:b2:anthropogenic-climate-change:budget}
\begin{enumerate}
  \item Trek de reeds bereikte menselijke opwarming, \qty{1.07}{\celsius} voor 2010--2019, af van de
        grens.
  \item Trek de opwarming af die de andere drijvende factoren zullen toevoegen tegen het moment dat de uitstoot van koolstofdioxide
        netto nul bereikt.
  \item Deel door de TCRE: dit is het budget vanaf het midden van 2010--2019.
  \item Trek de uitstoot sindsdien af, en een marge voor de trage terugkoppelingen van het aardsysteem die de
        modellen weglaten (het ontdooien van permafrost is de grootste).
\end{enumerate}
\end{method}

De beoordeling combineert op deze manier vijf componenten: de TCRE; de historische opwarming; een
extra opwarming door de andere drijvende factoren van ongeveer 0.1 tot \qty{0.2}{\celsius}; de opwarming die kan
aanhouden nadat de uitstoot stopt, gemiddeld op nul beoordeeld (\emph{waarschijnlijk} binnen $\pm
\qty{0.19}{\celsius}$); en de terugkoppelingen van het aardsysteem, ongeveer \qty{26}{Gt} koolstofdioxide per
graad. Voor \qty{1.5}{\celsius}, \qty{0.43}{\celsius} boven 2010--2019, is het resterende budget vanaf het
begin van 2020 \qty{500}{Gt} koolstofdioxide bij een waarschijnlijkheid van \qty{50}{\percent}, en
\qty{1350}{Gt} voor \qty{2}{\celsius}. De keuzes over de andere drijvende factoren verschuiven deze waarden met
\qty{220}{Gt} of meer; het weekendprobleem bouwt de eerste opnieuw op.
% ledger: ar6ch5:budget-parts, ar6spm:budget

\begin{dated}{2025}
% ledger: data:co2-mlo, data:gistemp (values read from the datasets as fetched 2026-10-22; offset computed)
Het jaargemiddelde van koolstofdioxide op Mauna Loa bereikte \qty{427.35}{ppm} in 2025, een forcering van
$3.93\log_2(427.35/278) = \qty{2.44}{W.m^{-2}}$ volgens de logaritmische wet. In de temperatuurreeks van NASA
is 2024 het warmste jaar, \qty{1.50}{\celsius} boven het gemiddelde van 1880--1900, de vroegste
jaren van die reeks; 2025 lag op \qty{1.40}{\celsius}.
\end{dated}

\begin{omfigure}
\centering
\begin{tikzpicture}
\begin{axis}[omrecord, width=0.86\linewidth, height=5.8cm, xmin=0, xmax=4000, ymin=0, ymax=2.6,
  xlabel={cumulatieve \ce{CO2}-uitstoot sinds 1850 (\unit{Gt})},
  ylabel={opwarming sinds 1850--1900 (\unit{\celsius})}, xtick={0,500,...,4000},
  /pgf/number format/1000 sep={\,}, xtick pos=bottom, ytick pos=left, clip=false,
  tick label style={font=\scriptsize}, label style={font=\small}]
% computed in figdata/bachelor-2/anthropogenic.py part tcre; ledger: ar6spm:tcre, ar6spm:budget
\addplot[name path=hi, draw=none, forget plot] table[x=E_GtCO2, y=dT_hi] {figdata/out/bachelor-2-anthropogenic-tcre.dat};
\addplot[name path=lo, draw=none, forget plot] table[x=E_GtCO2, y=dT_lo] {figdata/out/bachelor-2-anthropogenic-tcre.dat};
\addplot[omThm!18, forget plot] fill between[of=hi and lo];
\addplot[omThm, line width=1.3pt] table[x=E_GtCO2, y=dT_mid] {figdata/out/bachelor-2-anthropogenic-tcre.dat};
\draw[densely dashed, black!60] (axis cs:0,1.5) -- (axis cs:4000,1.5);
\draw[densely dashed, black!60] (axis cs:0,2.0) -- (axis cs:4000,2.0);
\node[anchor=north west, font=\scriptsize, text=black!70] at (axis cs:0,1.5) {\qty{1.5}{\celsius}};
\node[anchor=north west, font=\scriptsize, text=black!70] at (axis cs:0,2.0) {\qty{2}{\celsius}};
\fill[black] (axis cs:2390,1.07) circle (2.2pt);
\draw[black, line width=0.8pt] (axis cs:2150,1.07) -- (axis cs:2630,1.07);
\draw[black, line width=0.8pt] (axis cs:2390,0.8) -- (axis cs:2390,1.3);
\node[anchor=north west, font=\scriptsize, align=left] at (axis cs:2440,1.02) {1850--2019:\\\qty{2390}{Gt}, \qty{1.07}{\celsius}};
\draw[<->, omDef, thick] (axis cs:2390,1.62) -- (axis cs:2890,1.62);
\node[anchor=south, font=\scriptsize, text=omDef, align=center] at (axis cs:2640,1.64) {AR6: \qty{500}{Gt}};
\node[anchor=north west, font=\scriptsize, text=omThm, align=left] at (axis cs:150,2.5) {waarschijnlijk bereik van de TCRE,\\0.27 tot \qty{0.63}{\celsius} per \qty{1000}{Gt}};
\end{axis}
\end{tikzpicture}
\par\smallskip{\small Opwarming tegenover cumulatieve uitstoot van koolstofdioxide voor de beoordeelde TCRE,
\qty{0.45}{\celsius} per \qty{1000}{Gt} (lijn) en haar waarschijnlijke bereik (gearceerd), met de menselijke opwarming
en de uitstoot tot nu toe (punt, waarschijnlijke bereiken als balken). De pijl geeft het resterende budget aan voor
\qty{1.5}{\celsius} bij \qty{50}{\percent} vanaf 2020: kleiner dan een aflezing van de lijn alleen,
omdat de andere drijvende factoren ook opwarmen.}
\end{omfigure}


\section{Zeeniveau}\label{sec:b2:anthropogenic-climate-change:sea-level}

\begin{definition}[Zeeniveaubudget; thermosterische zeespiegelstijging]\label{def:b2:anthropogenic-climate-change:sea-level}
Het \emph{zeeniveaubudget}\index{zeeniveaubudget} is de vergelijking van de waargenomen verandering van het
globaal gemiddelde zeeniveau met de som van zijn geschatte bijdragen: de uitzetting van de opwarmende
oceaan, het afsmelten van gletsjers en van de Groenlandse en Antarctische ijskappen, en de verandering van het
op het land opgeslagen water. \emph{Thermosterische zeespiegelstijging}\index{thermosterische zeespiegelstijging} is het
deel dat te wijten is aan de thermische uitzetting van zeewater, bij constante massa.
\end{definition}

\begin{proposition}[Thermische uitzetting van een laag]\label{prop:b2:anthropogenic-climate-change:thermosteric}
Een oceaanlaag met dikte $H$ en thermische uitzettingscoëfficiënt $\alpha$, die bij constante massa en oppervlakte gelijkmatig
met $\Delta T$ opwarmt, stijgt met $\Delta h = \alpha\,\Delta T\,H$.
\end{proposition}
\begin{proof}
De massa $\rho H A$ van een kolom met oppervlakte $A$ is constant, dus $\Delta H/H = -\Delta\rho/\rho =
\alpha\,\Delta T$ in eerste orde, met $\alpha = -(1/\rho)\,\partial\rho/\partial T$.
\end{proof}
% ledger: teos10:alpha-beta

In de echte oceaan groeit $\alpha$ met temperatuur en druk, en de opwarming is geconcentreerd dicht bij
het oppervlak. De beoordeling vat het resultaat samen als een uitzettingsrendement van \qty{0.113}{m} per
\qty{e24}{J} warmte die aan de oceaan wordt toegevoegd, of \qty{0.617}{m} per graad opwarming van de hele oceaan
(\emph{middelgroot vertrouwen}). Over zijn gemiddelde diepte van \qty{3682}{m} is dat een effectieve $\alpha$ van
$0.617/3682 = \qty{1.68e-4}{\celsius^{-1}}$ (\cref{exo:b2:anthropogenic-climate-change:8}).
% ledger: ar6ch9:expansion, charette:ocean-volume

Het budget, getekend in de figuur hieronder, sluit binnen zijn onzekerheden. Over 1971--2018 tellen de bijdragen op tot
\qty{2.00}{mm.yr^{-1}}, tegenover een waargenomen stijging van \qty{2.33}{mm.yr^{-1}}; over 2006--2018 tellen ze
op tot 3.61 tegenover 3.69. Thermische uitzetting was de grootste afzonderlijke term, de helft van de som over 1971--2018, maar het ijs
haalt haar in: over 2006--2018 gaven gletsjers en ijskappen samen \qty{1.62}{mm.yr^{-1}}, meer dan
de uitzetting. De oceaan blijft ook in de diepte opwarmen lang nadat het oppervlak tot rust is gekomen, en de
ijskappen reageren over eeuwen; het zeeniveau zal onvermijdelijk eeuwen tot millennia blijven stijgen en zal
duizenden jaren hoog blijven (\emph{hoog vertrouwen}). Over de komende tweeduizend jaar wordt een
stijging geprojecteerd van ongeveer 2 tot \qty{3}{m} als de opwarming op \qty{1.5}{\celsius} blijft, 2 tot
\qty{6}{m} bij \qty{2}{\celsius} en 19 tot \qty{22}{m} bij \qty{5}{\celsius} (\emph{laag vertrouwen}).
% ledger: ar6ch9:table95, ar6spm:sea-commit

\begin{omfigure}
\centering
\begin{tikzpicture}
\begin{axis}[ybar stacked, width=0.72\linewidth, height=6.2cm, ymin=0, ymax=4.3, bar width=26pt,
  xtick=data, symbolic x coords={1971--2018, 1993--2018, 2006--2018}, enlarge x limits=0.25,
  ylabel={stijgsnelheid (\unit{mm.yr^{-1}})}, tick label style={font=\scriptsize},
  label style={font=\small}, ytick pos=left, xtick pos=bottom,
  legend style={at={(1.03,0.5)}, anchor=west, font=\scriptsize, draw=black!30},
  legend cell align=left, reverse legend]
% ledger: ar6ch9:table95
\addplot[fill=omDef!70, draw=omDef!90!black] coordinates {(1971--2018,1.01) (1993--2018,1.31) (2006--2018,1.39)};
\addlegendentry{thermische uitzetting}
\addplot[fill=omDef!30, draw=omDef!70!black] coordinates {(1971--2018,0.44) (1993--2018,0.55) (2006--2018,0.62)};
\addlegendentry{gletsjers}
\addplot[fill=omThm!55, draw=omThm!80!black] coordinates {(1971--2018,0.25) (1993--2018,0.43) (2006--2018,0.63)};
\addlegendentry{Groenlands ijs}
\addplot[fill=omThm!25, draw=omThm!70!black] coordinates {(1971--2018,0.14) (1993--2018,0.25) (2006--2018,0.37)};
\addlegendentry{Antarctisch ijs}
\addplot[fill=black!20, draw=black!50] coordinates {(1971--2018,0.15) (1993--2018,0.31) (2006--2018,0.60)};
\addlegendentry{landwater}
\node[fill=black, rotate=45, minimum size=5pt, inner sep=0pt] at (axis cs:1971--2018,2.33) {};
\node[fill=black, rotate=45, minimum size=5pt, inner sep=0pt] at (axis cs:1993--2018,3.25) {};
\node[fill=black, rotate=45, minimum size=5pt, inner sep=0pt] at (axis cs:2006--2018,3.69) {};
\end{axis}
\end{tikzpicture}
\par\smallskip{\small Het \omterm{def:b2:anthropogenic-climate-change:sea-level}{zeeniveaubudget}: waargenomen bijdragen aan de stijging van het globaal gemiddelde
zeeniveau, als snelheden over drie perioden die in 2018 eindigen, gestapeld, en de waargenomen snelheid (ruiten). Het
verschil tussen de stapel en de ruit is het deel dat het budget niet sluit. Waarden: IPCC Sixth
Assessment Report, Working Group~I, tabel 9.5.}
\end{omfigure}

\section*{Oefeningen}\label{sec:b2:anthropogenic-climate-change:exercises}

Gegevens voor de oefeningen, tenzij anders vermeld: forcering van verdubbelde koolstofdioxide $F_{2\times} =
\qty{3.93}{W.m^{-2}}$ vanaf $C_0 = \qty{278}{ppm}$; constante van Stefan--Boltzmann $\sigma =
\qty{5.670e-8}{W.m^{-2}.K^{-4}}$; straal van de aarde \qty{6371}{km}; één jaar
\qty{3.156e7}{s}; molaire massa's van koolstof en koolstofdioxide 12.011 en \qty{44.009}{g.mol^{-1}};
de terugkoppelingen en het \omterm{def:b2:anthropogenic-climate-change:sea-level}{zeeniveaubudget} van het hoofdstuk.
% ledger: ipcc:erf-2xco2, const:sigma, earth:Rmean, ciaaw:abridged-2024, ar6ch7:table710, ar6ch9:table95

\begin{exercise}[$\star$]\label{exo:b2:anthropogenic-climate-change:1}
De forcering van verdubbelde koolstofdioxide is \qty{3.75}{W.m^{-2}} zodra de stratosfeer zich heeft aangepast;
de troposfeer past zich daarna aan met $-0.60$ (temperatuur), +0.22 (waterdamp), +0.45 (wolken) en
+0.11 (landoppervlak) \unit{W.m^{-2}}. (a) Bepaal de ERF. (b) Met welk percentage overtreft ze de
\qty{3.71}{W.m^{-2}} van de vorige beoordeling? (c) Welke aanpassing is negatief, en waarom vermindert een warmere
troposfeer de forcering?
\end{exercise}
% ledger: ar6ch7:2xco2

\begin{exercise}[$\star$]\label{exo:b2:anthropogenic-climate-change:2}
Het klimaatsysteem won \qty{434.9}{ZJ} van 1971 tot 2018 en \qty{152.4}{ZJ} van 2006 tot 2018,
geteld tussen de middens van het eerste en het laatste jaar. (a) Bepaal de oppervlakte van de aarde.
(b) Zet elke winst om in een gemiddelde onbalans in \unit{W.m^{-2}}. (c) De oceaan nam
\qty{91}{\percent} op: hoeveel zettajoule over 1971--2018?
\end{exercise}
% ledger: ar6ch7:inventory, earth:Rmean

\begin{exercise}[$\star$]\label{exo:b2:anthropogenic-climate-change:3}
(a) Tel de centrale waarden van de terugkoppelingen in de figuur op, en vergelijk met de netto waarde. (b) Bepaal
de ECS met de netto waarde, en de opwarming met de \omterm{def:b2:anthropogenic-climate-change:feedbacks}{Planck-respons} alleen. (c) Met welke factor versterken
de andere terugkoppelingen de Planck-opwarming?
\end{exercise}

\begin{exercise}[$\star$]\label{exo:b2:anthropogenic-climate-change:4}
(a) Zet een TCRE van \qty{0.45}{\celsius} per \qty{1000}{Gt} koolstofdioxide om in graden per
\qty{1000}{PgC} koolstof, en doe hetzelfde voor haar waarschijnlijke bereik, 0.27 tot 0.63. (b) Zet de
\qty{2390}{Gt} koolstofdioxide die van 1850 tot 2019 werd uitgestoten, om in petagram koolstof. (c) Welke
opwarming geeft de TCRE voor deze uitstoot?
\end{exercise}
% ledger: ar6spm:tcre, ar6spm:budget

\begin{exercise}[$\star\star$]\label{exo:b2:anthropogenic-climate-change:5}
(a) Het oppervlak zendt \qty{398}{W.m^{-2}} uit en het broeikaseffect houdt daarvan \qty{159}{W.m^{-2}}
vast. Bepaal $g$ en de oppervlaktetemperatuur van een zwart lichaam dat \qty{398}{W.m^{-2}} uitzendt. (b) Bepaal de
\omterm{def:b2:anthropogenic-climate-change:feedbacks}{Planck-respons} bij \qty{288}{K}, en bij \qty{298}{K} met dezelfde $g$. (c) Zou de \omterm{def:b2:anthropogenic-climate-change:feedbacks}{Planck-respons}
alleen een warmere aarde meer of minder gevoelig maken?
\end{exercise}
% ledger: ipcc:budget-surface, ar6ch7:planck

\begin{exercise}[$\star\star$]\label{exo:b2:anthropogenic-climate-change:6}
Satellieten geven een \omterm{def:b2:anthropogenic-climate-change:feedbacks}{waterdampterugkoppeling} van $1.85 \pm 0.32$ en de studies een \omterm{def:b2:anthropogenic-climate-change:feedbacks}{terugkoppeling van de temperatuurgradiënt} van
$-0.50 \pm 0.20$ \unit{W.m^{-2}.\celsius^{-1}} (één standaardafwijking). (a) Bepaal hun som.
(b) Als de twee onafhankelijk waren, wat zou dan de standaardafwijking van de som zijn? (c) De modellen
geven de som een standaardafwijking van slechts \num{0.10}. Wat zegt dat over de twee fouten?
\end{exercise}
% ledger: ar6ch7:wv-lr

\begin{exercise}[$\star\star$]\label{exo:b2:anthropogenic-climate-change:7}
Van 1850--1900 tot 2006--2019 is $\Delta N = \qty{0.59}{W.m^{-2}}$, $\Delta T = \qty{1.03}{\celsius}$
en $\Delta F = \qty{2.20}{W.m^{-2}}$. (a) Bepaal $\lambda$ en de effectieve gevoeligheid. (b) Verwacht wordt dat de
terugkoppelingsparameter met $\lambda' = \qty{0.5}{W.m^{-2}.\celsius^{-1}}$ minder negatief wordt
naarmate het opwarmingspatroon verandert. Bepaal de gevoeligheid $-F_{2\times}/(\lambda + \lambda')$.
(c) Herhaal (b) met $\lambda' = 0.1$.
\end{exercise}
% ledger: ar6ch7:energy-budget, ar6ch7:pattern, ar6ch7:pattern-ecs

\begin{exercise}[$\star\star$]\label{exo:b2:anthropogenic-climate-change:8}
(a) Bepaal uit de \qty{0.617}{m} stijging per graad opwarming van de hele oceaan, met gemiddelde diepte
\qty{3682}{m}, de effectieve uitzettingscoëfficiënt. (b) De oceaan won \qty{396}{ZJ} over
1971--2018. Bepaal met een uitzettingsrendement van \qty{0.113}{m} per \qty{e24}{J} de thermosterische
stijging, en vergelijk met de \qty{47.5}{mm} van het budget. (c) Een laag van \qty{700}{m} warmt
\qty{0.2}{\celsius} op met $\alpha = \qty{2e-4}{\celsius^{-1}}$: hoeveel verhoogt ze het zeeniveau?
\end{exercise}
% ledger: ar6ch9:expansion, charette:ocean-volume, ar6ch7:inventory, ar6ch9:table95

\begin{exercise}[$\star\star$]\label{exo:b2:anthropogenic-climate-change:9}
Een modelvingerafdruk over vier decennia is $x = (0.10, 0.20, 0.30, 0.40)$ \unit{\celsius}; de
waarnemingen zijn $y = (0.12, 0.18, 0.33, 0.41)$ \unit{\celsius}. (a) Bepaal de
kleinstekwadratenschaalfactor $\beta$ van $y = \beta x + \varepsilon$. (b) Zijn standaardfout, geschat uit ongeforceerde
simulaties, is \num{0.08}. Is de respons gedetecteerd? Is ze in overeenstemming met het model? (c) Wat
zou een $\beta$ van 0.5 met dezelfde fout betekenen?
\end{exercise}

\begin{exercise}[$\star\star\star$]\label{exo:b2:anthropogenic-climate-change:10}
Twee vingerafdrukken, broeikasgassen $g = (0.2, 0.4, 0.6, 0.8)$ en aerosolen $a = (-0.1, -0.2,
-0.2, -0.3)$, beide in \unit{\celsius} over vier decennia, worden aangepast aan de waarnemingen $y = (0.1,
0.2, 0.4, 0.5)$. (a) Schrijf de normaalvergelijkingen van de kleinstekwadratenaanpassing $y = \beta_g g + \beta_a a
+ \varepsilon$ en los ze op. (b) Bepaal de opwarming die in het laatste decennium aan elke factor wordt toegeschreven.
(c) Waarom groeien de onzekerheden op $\beta_g$ en $\beta_a$ wanneer de twee vingerafdrukken op elkaar lijken?
\end{exercise}

\begin{exercise}[$\star\star\star$]\label{exo:b2:anthropogenic-climate-change:11}
Gebruik de methode van het \omterm{def:b2:anthropogenic-climate-change:budget}{resterende koolstofbudget} met de TCRE van \qty{0.45}{\celsius} per
\qty{1000}{Gt}, de \qty{210}{Gt} die van 2015 tot 2019 werd uitgestoten en \qty{26}{Gt} per graad voor de
terugkoppelingen van het aardsysteem. (a) Bepaal het budget voor \qty{2}{\celsius}, \qty{0.93}{\celsius} boven
2010--2019, als de andere factoren \qty{0.1}{\celsius} toevoegen. (b) Welke extra opwarming door de andere
factoren maakt het gelijk aan de beoordeelde \qty{1350}{Gt}? Ligt dat binnen het beoordeelde bereik? (c) Toon aan
dat elke \qty{0.1}{\celsius} opwarming door andere factoren dan CO$_2$ elk budget met hetzelfde bedrag verandert.
\end{exercise}

\begin{exercise}[$\star\star\star$]\label{exo:b2:anthropogenic-climate-change:12}
De netto terugkoppeling heeft een zeer waarschijnlijk bereik van $-1.81$ tot \qty{-0.51}{W.m^{-2}.\celsius^{-1}}, rond
$-1.16$. (a) Bepaal de ECS aan de twee uiteinden en in het midden. (b) Toon aan dat een bereik van $\lambda$
dat symmetrisch is rond zijn midden een ECS-bereik geeft dat dat niet is, en zeg welke kant langer is. (c) Waarom
vinden de beoordelingen het moeilijker een hoge ECS uit te sluiten dan een lage?
\end{exercise}
% ledger: ar6ch7:table710

\begin{problem}[{Weekendprobleem --- Het budget lezen}]\label{pb:b2:anthropogenic-climate-change:1}
Dit probleem volgt de energie en de koolstof van de beoordeling van de forcering tot het resterende
budget voor \qty{1.5}{\celsius}. Gebruik de gegevens van de oefeningen, en: ERF 1750--2019 (\unit{W.m^{-2}})
koolstofdioxide 2.16, methaan 0.54, distikstofoxide 0.21, gehalogeneerde gassen 0.41, ozon 0.47,
stratosferische waterdamp 0.05, aerosol--straling $-0.22$, aerosol--wolken $-0.84$, landgebruik
$-0.20$, roet op sneeuw 0.08, condensstrepen 0.06; koolstofdioxide \qty{410}{ppm} in 2019 en
\qty{427.35}{ppm} in 2025; warmteopname \qty{0.79}{W.m^{-2}} over 2006--2018; de veranderingen van de energiebalans
$\Delta N = 0.59$, $\Delta F = \qty{2.20}{W.m^{-2}}$, $\Delta T = \qty{1.03}{\celsius}$; TCRE
\qty{0.45}{\celsius} per \qty{1000}{Gt} koolstofdioxide (waarschijnlijk 0.27 tot 0.63); \qty{2390}{Gt}
uitgestoten 1850--2019; menselijke opwarming \qty{1.07}{\celsius} voor 2010--2019; \qty{210}{Gt} uitgestoten van
2015 tot 2019; extra opwarming door andere factoren dan CO$_2$ \qty{0.1}{\celsius}, de onderkant van het beoordeelde bereik 0.1
tot \qty{0.2}{\celsius}; terugkoppelingen van het aardsysteem \qty{26}{Gt} per graad. Rond af op drie beduidende
cijfers.
% ledger: ipcc:erf-agents, ipcc:erf-2xco2, data:co2-mlo, ar6ch7:inventory, ar6ch7:energy-budget, ar6spm:tcre, ar6spm:budget, ar6ch5:budget-parts

\noindent\textbf{Deel I --- De forcering.}
\begin{enumerate}
  \item Tel de forcering van de vier goed gemengde broeikasgassen op.
  \item Tel de forcering van de aerosolen op.
  \item Tel alle menselijke termen op. Vergelijk met de beoordeelde \qty{2.72}{W.m^{-2}}.
  \item Controleer de forcering van koolstofdioxide in 2019 met de logaritmische wet.
  \item Bepaal de forcering van koolstofdioxide in 2025 met dezelfde wet.
  \item Welke term heeft de grootste onzekerheid, en waarom overheerst die de onzekerheid van het totaal?
\end{enumerate}

\noindent\textbf{Deel II --- De energiebalans.}
\begin{enumerate}[resume]
  \item Hoeveel energie won het klimaatsysteem over 2006--2018, in zettajoule?
  \item Hoeveel daarvan ging naar de oceaan?
  \item Bepaal de terugkoppelingsparameter uit de veranderingen van de energiebalans.
  \item Bepaal de effectieve \omterm{def:b2:anthropogenic-climate-change:ecs}{evenwichtsklimaatgevoeligheid}.
  \item Bepaal de \omterm{def:b2:anthropogenic-climate-change:tcr}{transiënte klimaatrespons} uit dezelfde gegevens.
  \item Bepaal de ECS uit de netto terugkoppeling van de figuur, en vergelijk de twee waarden.
  \item Geef de fysische reden waarom de historische schatting de laagste is.
\end{enumerate}

\noindent\textbf{Deel III --- Opwarming per ton.}
\begin{enumerate}[resume]
  \item Zet de TCRE om in graden per \qty{1000}{PgC} koolstof.
  \item Zet haar waarschijnlijke bereik om.
  \item Zet de uitstoot van 1850--2019 om in petagram koolstof.
  \item Welke opwarming geeft de TCRE voor de uitstoot van 1850--2019?
  \item Vergelijk met de menselijke opwarming van \qty{1.07}{\celsius}. Wat zegt de vergelijking over het
        netto effect van de andere menselijke factoren?
\end{enumerate}

\noindent\textbf{Deel IV --- Het resterende budget voor \qty{1.5}{\celsius}.}
\begin{enumerate}[resume]
  \item Welke opwarming blijft er over boven 2010--2019?
  \item Trek de opwarming door andere factoren dan CO$_2$ af en deel door de TCRE: welk budget blijft er over vanaf het midden
        van 2010--2019?
  \item Trek de uitstoot van 2015--2019 af.
  \item Trek de terugkoppelingen van het aardsysteem over de resterende opwarming af.
  \item Hoeveel verandert het budget als de andere factoren in plaats daarvan \qty{0.2}{\celsius} toevoegen?
  \item Samengevat: wat is het \omterm{def:b2:anthropogenic-climate-change:budget}{resterende koolstofbudget} voor \qty{1.5}{\celsius}, vanaf het begin van
        2020, berekend uit de TCRE? Vergelijk met de beoordeelde \qty{500}{Gt}.
\end{enumerate}
\end{problem}
